\section{The reducible characteristic-\(\ell\) curves require no new cubic finiteness theorem}
\label{sec:cubic-typeii-curve-exclusion}

Retain the exact old datum \(D\), coefficient DVR \(O\),
marked elements \(c,g_1\), and cubic normalized ring \(S\)
of AC-G49--AC-G59.  Write
\[
 \Gamma=G_{\mathbb Q},\qquad H=G_K=\ker\eta,
 \qquad \gamma=D.\Psi,\qquad
 q(hc^a)=\chi(h),\qquad j\in\{1,2\}.
 \tag{G60.1}
\]
Here \(K=\mathbb Q(\omega)\), \(O\) contains \(\mu_3\),
\(\ell\ne2,3,269\), and all original \texttt{GD} and local
published-theorem conditions remain in force.
The old datum is formed over this same coefficient ring.
In addition, retain the exact source assertion
\[
 \operatorname{FinEnd}(\ell,\eta,\gamma,c,g_1,B_\ell)
 \tag{G60.2}
\]
for these data.  This is the assertion in
\path{CurveDVRSpecialization/RatTorsionFree.lean},
\texttt{UpperTheory.V.Pr.FinEnd}; its quantifiers are
recalled below.  In particular its multiplier is the
actual \(D.\eta\), not an implicitly substituted
\(\eta_4\).  Its arithmetic supply and all hypotheses of
the supplied assembly are retained.
Naming the source proposition is not a proof of its
arithmetic supply: (G60.2) is an explicit retained input.
Neither its \texttt{AutEnd3} nor its \texttt{AutEnd2}
companion is used in the argument below.

Let \(f:S\to V\) be an actual characteristic-\(\ell\)
complete-DVR curve as in (G59.2), and put \(r=f_*r_S\).
Its residual characters on \(H\) are
\(\bar\chi^{\,j}\bar\gamma,1,
\bar\chi^{\,j}\bar\gamma^{-1}\).
An element of \(O\) inside a formula over \(V\) means its
image under the given local algebra map.
All source paths in this section have the pinned root
\path{Development/AutomorphicArithmetic/AutomorphicVarieties/}.

\begin{proposition}[AC-G60: endpoint untwisting of characteristic-\(\ell\) block curves]
\label{AC-G60}
Suppose the actual curve \(r\) has the endpoint/middle
block decomposition on \(H\), with no hypothesis that
its middle character is already trivial.
Then
\[
 T(g)=\operatorname{diag}(q(g)^{-j},1,q(g)^{-j}),\qquad
 \widetilde r(g).\mathrm{mat}=T(g)r(g).\mathrm{mat},\qquad
 \widetilde r(g).\mathrm{mult}=\eta(g)
 \tag{G60.3}
\]
defines an ordinary curve over the test \(O\)-algebra
\(V\) in the exact sense \texttt{OrdPt} for the old datum.
Consequently every such block curve has the full seed
characteristic polynomial
\[
 \det(X-r(h).\mathrm{mat})
  =(X-\bar\chi(h)^j\bar\gamma(h))(X-1)
    (X-\bar\chi(h)^j\bar\gamma(h)^{-1})
       \quad(h\in H).
 \tag{G60.4}
\]
In particular no actual characteristic-\(\ell\) curve can
have Type II shape over an algebraic closure of
\(\operatorname{Frac}V\).
\end{proposition}
\begin{proof}
\proofstep{The homomorphism and the literal residual fields.}
The block condition says that all entries between
\(V e_1\) and \(V e_0\oplus V e_2\) vanish.
It holds on all of \(\Gamma\): if \(g=hc\), then
right multiplication by \(c\), whose matrix is \(I\),
gives \(r(hc).\mathrm{mat}=r(h).\mathrm{mat}\).
Thus \(T(g')\) commutes with every \(r(g).\mathrm{mat}\).
For the source involution \(\theta(A)=JA^{-t}J\),
\(\theta(T(g'))=T(g')^{-1}\).
The anticyclotomic crossed identity
\(q(gg')=q(g)q(g')^{\eta(g)}\) therefore gives
\[
 \begin{split}
 T(gg')r(gg').\mathrm{mat}
  ={}&T(g)r(g).\mathrm{mat}\,
       \theta_{\eta(g)}
        \bigl(T(g')r(g').\mathrm{mat}\bigr).
 \end{split}
 \tag{G60.5}
\]
This proves the homomorphism assertion over \(V\) itself.
Continuity is preserved because \(q\) is locally constant
and its values are units.  Multiplication by \(T\) cancels
the cubic factors in every residual endpoint entry and
leaves the multiplier and middle entry unchanged.
Hence \texttt{Red} is exactly the old \(r_0(\eta,\gamma)\)
residual condition, over the same test algebra and residue
map.  No change of coefficient characteristic is involved.

The value at \(c\) is still \(I\).
The two mixed entries of
\(\widetilde r(g_1).\mathrm{mat}
 -\widetilde r(cg_1^{-1}c).\mathrm{mat}\)
vanish because both matrices are block diagonal.
This is exactly \texttt{Normalised}; it imposes no new
condition on \(\chi(g_1)\).

\proofstep{The exact away-inertia condition.}
The source \texttt{OrdPt.away} asks for identity at every
\(g\in I_v\), \(v\nmid\ell\), such that
\(\eta(g)=1\) and \(\gamma(g)=1\).
For such a \(g\), the cubic residual matrix has endpoints
\(\bar q(g)^j\) and middle \(1\); its cube is \(I\).
Thus \(g^3\) belongs to the actual modified killing set
\(\mathcal K_\chi\) of (G49.2), and \(r(g)^3=1\).
Since \(T(g)^3=I\) and \(T(g)\) commutes with the block
matrix \(r(g).\mathrm{mat}\), the matrix
\(C=\widetilde r(g).\mathrm{mat}\) satisfies
\[
        C^3=I,\qquad C\equiv I\pmod{\mathfrak m_V}.
 \tag{G60.6}
\]
The matrix \(C^2+C+I\) reduces to \(3I\).
Because \(\ell\ne3\), its determinant is a unit in
\(V\), and the matrix is invertible.
The identity \((C-I)(C^2+C+I)=0\) therefore gives
\(C=I\).  The multiplier is already \(1\), so
\(\widetilde r(g)=1\).
This verifies every quantified place and inertia element
in \texttt{OrdPt.away}.  It does not assume that
\(\psi\) is unramified at \(269\), that the middle
character is trivial, or that \(\gamma\) has any new
order restriction.

\proofstep{The strict integral ordinary flag.}
Use the explicit choice of \(s\in I_\ell\) and
projectors \(E_i\) from (G57I.1)--(G57I.3).
Every \(E_i\) is a polynomial in the block matrix
\(r(s).\mathrm{mat}\).  Hence the strict basis matrix
\[
                  h=(E_0e_0,E_1e_1,E_2e_2)
 \tag{G60.7}
\]
is block diagonal for the same endpoint/middle
decomposition.  It is invertible and congruent to \(I\).
By AC-G57-integral it makes \(r\) upper triangular on the
local \(\ker\eta\) and gives the ordered inertia diagonal
\(\gamma,1,\gamma^{-1}\).
The matrix \(h\) commutes with every \(T(g)\).
On the local \(\ker\eta\), semidirect conjugation is
ordinary matrix conjugation, so
\[
 h^{-1}\widetilde r(g).\mathrm{mat}\,h
     =T(g)\,h^{-1}r(g).\mathrm{mat}\,h.
 \tag{G60.8}
\]
The right side is upper triangular.
Moreover \(q=1\) on \(I_\ell\): the cubic character is
unramified at the retained good coefficient primes,
and \(\eta\) is unramified there.
Thus the ordered diagonal on inertia is unchanged.
Together with \(h\equiv I\), this is precisely
\texttt{InHat} and \texttt{TriIn}, hence \texttt{CondT},
in
\path{UnitaryThreeDeformations/AdjointDiagonalRepPreservesWeightPiece.lean}.
We have now checked all five fields of the literal
\texttt{OrdPt}: residual condition, continuity,
normalization, away-inertia killing, and strict local
ordinary condition.
This argument needs no membership in the old selected
ideal \(I_t\), and asserts no such membership.

\proofstep{Apply the unchanged old finiteness assertion.}
The quantified test algebra in (G60.2) is any
\texttt{IsCD} \(O\)-algebra annihilating
\(\mathfrak m_O\); this is exactly \(V\).
It supplies a finite-index subgroup \(U\subset\Gamma\)
on which the characteristic polynomial of
\(\widetilde r\), on \(H\), equals that of the old seed.
Intersect with \(\ker(\chi)\subset H\), a finite-index
subgroup of \(\Gamma\).  On
\(U'=U\cap\ker(\chi)\), untwisting is the identity
and the old and cubic seeds coincide.  Therefore
\[
 \det(X-r(u).\mathrm{mat})
        =\det(X-\bar r_\chi(u).\mathrm{mat})
                     \qquad(u\in U').
 \tag{G60.9}
\]
For each \(h\in H\), some positive power \(h^n\)
belongs to \(U'\).  Every eigenvalue of
\(r(h).\mathrm{mat}\) has its \(n\)-th power among
roots of the right side of (G60.9), so it is algebraic
over the finite field \(k\).  This implication is
valid even when \(\ell\mid n\).
The characteristic coefficients are consequently
algebraic over \(k\).  Each belongs to \(V\) and has
its prescribed residue in \(k\), so it equals that
residue by the finite-field argument in
AC-G59-polynomials.  This proves (G60.4).

If the curve had Type II shape, its middle character,
being a root of (G60.4) with residue \(1\), would equal
\(1\).  Canceling \(X-1\) would give its irreducible
endpoint block the same characteristic polynomials
as the two displayed seed characters.
The elementary image-algebra trace argument proved in
AC-G59-polynomials excludes this in dimension two.
Thus the Type II case is impossible.
\end{proof}

\begin{corollary}[AC-G60-triangular: the diagonal of every triangular curve is fixed]
\label{AC-G60-triangular}
Under the same exact old \texttt{FinEnd} hypothesis,
every actual upper or lower triangular
characteristic-\(\ell\) curve has the cubic seed diagonal.
\end{corollary}
\begin{proof}
Suppose first that \(r\) is upper triangular on \(H\).
Right multiplication by \(c\) makes every matrix on
\(\Gamma\) upper triangular as well.
The involution \(\theta(A)=JA^{-t}J\) preserves upper
triangular matrices and acts on their diagonal by
reversed inversion.  Consequently taking the diagonal
is a homomorphism from this triangular semidirect group
to its diagonal semidirect subgroup.
Define \(r^{\rm diag}\) by this operation.
It retains the exact multiplier, residual matrix,
continuity, normalization, and killing of
\(\mathcal K_\chi\).

The matrix \(h\) in (G57I.3) is now upper triangular:
each projector is a polynomial in an upper triangular
matrix, and its \(i\)-th column has no entries below
row \(i\).  Conjugation by an invertible upper
triangular matrix preserves each diagonal entry of
an upper triangular matrix.
AC-G57-integral therefore says that the diagonal of
\(r(s)\), for every \(s\in I_\ell\), already is
\(\gamma(s),1,\gamma(s)^{-1}\) in the original basis.
Thus \(r^{\rm diag}\) has the exact local ordinary
condition with strict basis \(I\).
The same argument for lower triangular matrices uses
lower triangular projectors and preserves the same
ordered diagonal.

Apply the proof of AC-G60 to this diagonal representation.
That proof used only its displayed residual,
normalization, killing and block-compatible local
ordinary properties, all just verified; it did not use
an additional factorization through \(S\).
Its characteristic polynomial is therefore the full
cubic seed polynomial at every \(h\in H\).
Each diagonal entry is algebraic over \(k\) and has
its prescribed residual coordinate value, so equals
that value.  These are the diagonal entries of the
original curve.  Right multiplication by \(c\)
extends the assertion to the other coset when needed.
\end{proof}

\begin{corollary}[AC-G60-M2: only the Type I characteristic-\(\ell\) curves remain]
\label{AC-G60-M2}
Under (G60.2), the actual condition \texttt{M2} is
equivalent to the nonexistence of an actual
characteristic-\(\ell\) curve whose restriction to
\(H\), after scalar extension to an algebraic closure
of its fraction field, is irreducible of dimension three.
Thus the new cubic arithmetic input needed for this
part of the construction is confined to Type I curves.
\end{corollary}
\begin{proof}
If \texttt{M2} holds, AC-G59 makes every actual curve
triangular, excluding Type I.
Conversely assume no Type I actual curve exists.
The seed-independent \texttt{shape\_cases} theorem,
with the checked normalization and marked nonzero
\texttt{zeta}, gives upper, lower, or Type II shape
for each curve after algebraic scalar extension.
The Type II option is excluded by AC-G60: its block
entries vanish over the extension and hence already
in the domain \(V\).
The upper or lower entries similarly vanish already
over \(V\), and AC-G60-triangular supplies the seed
diagonal.  The exact integral criterion AC-G59 now
gives \texttt{M2}, including its vertical torsion.
\end{proof}

\paragraph{The exact arithmetic supply and the remaining gate.}
\path{CurveDVRSpecialization/RatTorsionFree.lean}
defines \texttt{OrdPt} and \texttt{FinEnd} with arbitrary
\(\eta\).  Its \texttt{HostEnd} clause specifically uses
\(\eta_4\), and must not be substituted for (G60.2).
The separate \texttt{GdEnd} clause, and
\texttt{GdEndAlt} in
\path{ModularityArgument/ThresholdArguments/DenseThresholdCases/AutEnd3AltAlt.lean},
quantify over the actual \(D.\eta\) and \(D.\Psi\).
The complete assembly proof
\path{ModularityArgument/ThresholdArguments/DenseThresholdCases/GdEndOf.lean},
\texttt{gdEnd\_of\_frame\_linkedEnd\_downCsdK\_descent},
constructs that \texttt{FinEnd} from its retained
\texttt{GdFrameEx} and \texttt{LinkedEndEAlt} data through
\texttt{finEnd\_of\_finPt}.
All of those arithmetic premises stay in force;
no claim is made that hypothesis A alone supplies them.
The marked element must be the one for which the old
assertion holds, which is permitted when constructing
the actual cubic ring.

The new reduction concerns characteristic-\(\ell\)
curves and integral degree-zero finiteness.
It does not prove the absence of Type I curves,
characteristic-zero Type I automorphy, the mixed
Type II tangent vanishing, or a numerical Selmer bound.
In particular the old finite-index conclusion is
used only after an actual ordinary representation
has been constructed in its exact quantified domain.

\paragraph{Geometric dictionary.}
Endpoint untwisting is tensoring the endpoint lisse
local system by the inverse finite cubic character,
while retaining the middle summand.  This uses the
already audited SGA~1, Expos\'e~V, \S\S4 and~7
representation and finite-cover dictionary.
The strict integral flag is the polynomial-projector
construction of AC-G57-integral.  The reduction of
integral support to all characteristic-\(\ell\)
DVR curves is the exact source curve theorem of
AC-G59, not a consequence of generic-fiber dimension
or of the categorical dictionary alone.
